\subsection*{一步一步建立直观}

\paragraph{1. 先固定数据来源，再评价一种模型}

真实分布 $p$ 决定结果出现的频率，模型分布 $q$ 决定各结果被赋予的理想码长。交叉熵因此是按 $p$ 求平均、对 $q$ 取负对数。右图让两者分别可调，避免把“模型认为罕见”和“现实中确实罕见”混为一谈。



\paragraph{2. 错把常见结果当罕见，会付出更多代价}

保持真实概率不变，拖动模型点偏离它。某些常出现的结果被分配到更低概率，负对数损失随之变大。蓝线高出熵基线的部分是 KL 散度；在正概率条件下，它非负，模型等于真实分布时为零。这是平均陈述，不保证每个单独样本损失都最小。



\paragraph{3. 更低的最优损失不总说明模型更强}

把真实分布改为偏斜分布，并让模型准确匹配。最优交叉熵随真实熵一起下降，因为数据更容易预测，而非模型结构发生了进步。比较不同任务的损失时，应考虑数据分布、单位与分词方式。窗口只验证二元分布关系，不直接评价大型语言模型的能力。



\subsection*{定义、命题与证明}

\paragraph{定义：有限分布的交叉熵与 KL 散度}

对正概率分布 $p,q$，定义 $H(p,q)=-\sum_i p_i\log_2q_i$，$H(p)=-\sum_i p_i\log_2p_i$，$D(p\Vert q)=\sum_i p_i\log_2(p_i/q_i)$。

\paragraph{命题：模型失配增加平均编码代价}

有 $H(p,q)=H(p)+D(p\Vert q)$，且 $D(p\Vert q)\ge0$；在此正概率情形，等号当且仅当 $p=q$。

\begin{proof}

把 $\log_2(p_i/q_i)$ 拆成两个对数并求和，立即得到分解。为证明非负，用 $\ln u\le u-1$，代入 $u=q_i/p_i$，得到 $p_i\ln(p_i/q_i)\ge p_i-q_i$。

求和后右侧为零，再除以正数 $\ln2$ 即得结论。对数不等式仅在 $u=1$ 时取等，因此所有正权重项同时取等恰好要求 $p_i=q_i$。含零概率时需另外约定边界与无穷代价。

\end{proof}
